% Source: https://artofproblemsolving.com/community/c3h2754539_3_problems
% Stable identity: aops-topic-2754539-q01
On the line, $y=-\dfrac x2$. Substituting, the $x$-coordinates of $A$ and $B$ are the roots $x_1,x_2$ of
\[x^2+\dfrac x2-a=0,\]
so $x_1+x_2=-\dfrac12$ and $x_1x_2=-a$. Then $(x_1-x_2)^2=\dfrac14+4a$.

\begin{asy}
usepackage("amsmath");
size(6cm);
defaultpen(fontsize(10pt));
import graph;
real a=11/80;
real f(real x){return x^2-a;}
draw(graph(f,-0.85,0.75),linewidth(0.8));
draw((-0.9,0.45)--(0.8,-0.4));
draw((-0.95,0)--(0.85,0),Arrow(4));draw((0,-0.25)--(0,0.65),Arrow(4));
real x1=(-0.5-sqrt(0.25+4a))/2,x2=(-0.5+sqrt(0.25+4a))/2;
pair A=(x1,-x1/2),B=(x2,-x2/2);dot(A);dot(B);
label("$A$",A,SW,fontsize(8pt));label("$B$",B,S,fontsize(8pt));label("$y=x^2-a$",(0.75,f(0.75)),E,fontsize(7pt));label("$x+2y=0$",(0.8,-0.4),E,fontsize(7pt));
\end{asy}

Since both points are on a line of slope $-\dfrac12$,
\[AB^2=(x_1-x_2)^2\left(1+\dfrac14\right)=\dfrac54\left(\dfrac14+4a\right)=1.\]
So $\dfrac14+4a=\dfrac45$ and $a=\boxed{\dfrac{11}{80}}$.
